\chapterrule
\section{外积与矩阵：如何真正“写入”一条记忆？}
\label{intro:04-key-value-outer-product}
\chaptercompass{一个矩阵怎样保存 Key 和 Value 的关联？}{外积写入 $S=k^\top v$，行查询从左边读取 $xS$。}{$k:1\times d_k,\ v:1\times d_v\to S:d_k\times d_v\to xS$}{这里写入的是示例状态矩阵；是否训练取决于应用场景}
\modelpipelinefocus{3}{Attention 的记忆视角}
\phantomsection\label{concept:outer-product}

Transformer 主线里每个 Token 占一行。现在换一种视角理解线性层：
能否把“遇见 Key 就读出 Value”的关联写入一个矩阵？
和前面一样，Key、Query 与 Value 都以\textbf{一条记录占一行}来表示。

\subsection{先问输入、输出和形状}
单次 Query 是 $x:1\times d_k$，Key 为 $k:1\times d_k$，
Value 为 $v:1\times d_v$。我们希望读出一行 $1\times d_v$ 的结果。
\[
\underbrace{x}_{1\times d_k}
\underbrace{k^\top}_{d_k\times1}
\underbrace{v}_{1\times d_v}
=(xk^\top)v.
\]
先用点积得到匹配分数，再乘 Value；不能直接把 Key 和 Value 的维度混同。

\subsection{亲手写入一条二维记忆}
\localfixture{独立二维行向量例}
令 $k=(1,0)$，$v=(2,3)$，$x=k$。
两条向量的外积矩阵是
\[
S=k^\top v
=\begin{bmatrix}1\\0\end{bmatrix}
\begin{bmatrix}2&3\end{bmatrix}
=\boxed{\begin{bmatrix}2&3\\0&0\end{bmatrix}}.
\]
读取时从左边乘一条\textbf{行向量}：
\begin{equation}
\boxed{xS=(1,0)S=(2,3).}
\end{equation}
因为 $kk^\top=1$，所以 $kS=(kk^\top)v=v$；
若 Key 并非单位向量，结果一般会乘上 $\|k\|^2$。
外积得到一个矩阵，点积得到一个数；二者不要混淆。

\begin{idea}
这就是同一计算的两种写法：$x(k^\top v)=(xk^\top)v$。
如果有 $M$ 条记忆，Key 和 Value 分别\textbf{按行}存进
$K:M\times d_k$、$V:M\times d_v$，
匹配权重为 $\vect a^\top:1\times M$，
则最后读出为 $\vect a^\top V:1\times d_v$。
形式直接对齐主线 Attention 的“权重乘 Value”。
\end{idea}

\chapterchecktasks{04-key-value-outer-product}{形状／手算／预测／找错}{
\item \textbf{形状：}行向量 Key $k:1\times2$、Value $v:1\times3$。外积写入 $S=k^\top v$ 的形状是什么？查询 $q:1\times2$ 后 $qS$ 又是什么形状？
\item \textbf{手算：}$k=(1,0),v=(0,1)$，写出矩阵 $S=k^\top v$，再用 $q=(1,0)$ 读取。
\item \textbf{预测：}同一条记忆若把 Key 写入时改为 $k'=(0.5,0)$，再以 $q=k'$ 读取，恢复的 Value 是多少倍？
\item \textbf{找错：}同学把 $S$ 写成 $vk^\top$，仍称它是按行向量约定可供 $qS$ 查询的状态。说明为何形状或方向不匹配。
}
\chapterrecallanswer{行向量 Key 的外积写入方向是 $k^\top v$；查询 $qS$ 得到 Value 空间的一行。}


